stručný popis každého modulu (informačně obohacující věta na dvojici .c/.h)

%DAmian%
\section{hashtable.c/.h}
Modul pre prácu s hash tabuľkou, ktorá slúži na ukladanie a vyhľadávanie dát.

\section{codegen.c/.h}
Modul pre generovanie finálneho kódu z syntaktického stromu.

\section{gen\_func.c/.h}
Modul pre generovanie jednotlivých častí kódu vo finálnom prekladači.

\section{codegen\_helper.c/.h}
Modul s pomocnými funkciami pre generovanie kódu.

\section{helper\_func.c/.h}
Modul s pomocnými funkciami, ktoré sa používajú v rôznych častiach projektu.

%Daniel%
\section{symtable.c/ .h}
Tabuľka symbolov ako zásobník hashovacích tabuliek, ktorý ukladá informácie o deklarovaných premenných a funkciách, spravuje rozsahy ich platnosti. Pri ukladaní premenných vytvára ich unikátne manglované názvy (GF, LF).

\section{semantic.c/ .h}
Rekurzívne prechádza strom AST, pomocou tabuľky symbolov overuje deklarácie, kompatibilitu operácií, pravidlá volaní... Dopĺňa do stromu uzly pre implicitné konverzie a dopĺňa odvodené dátové typy.

%Michal%
\section{parser.c/ .h}
Implementuje syntaktickú analýzu zhora nadol metódou rekurzívneho zostupu s vyrovnávacou pamäťou pre jeden dopredný token (lookahead). Kontroluje syntaktickú štruktúru programu podľa LL gramatiky, riadi prechod do precedenčnej analýzy pri spracovaní výrazov a postupne vytvára abstraktný syntaktický strom (AST).

\section{expression.c/ .h}
Zabezpečuje syntaktickú analýzu výrazov zdola nahor pomocou precedenčnej syntaktickej analýzy riadenej vlastným zásobníkom (PrecStack) a precedenčnou tabuľkou operátorov (precedence\_table). Spracúva unárne aj binárne operátory, uzátvorkované podvýrazy a vnorené volania funkcií, pričom pri redukciách vytvára príslušné podstromy AST.

\section{ast.c/ .h}
Definuje dátové štruktúry (ASTNode, ASTNodeType) a rozhranie na vytváranie uzlov AST reprezentujúceho prekladaný program. Súrodenecké uzly (príkazy v bloku, parametre, argumenty a definície funkcií) prepája do obojsmerne viazaného zoznamu (next, prev), zatiaľ čo všetky alokované uzly zároveň eviduje v intruzívnom jednosmerne viazanom zozname (alloc\_next) na ich centralizované uvoľnenie vrátane alokovaných reťazcov pri ukončení prekladu alebo pri chybe.

\newpage